% ============================================================================
% Conclusão
% ============================================================================
\newpage
\section*{CONCLUSÃO}
\addcontentsline{toc}{section}{CONCLUSÃO}

Este trabalho teve como objetivo desenvolver uma solução tecnológica corporativa
integrada para apoiar a transformação digital de uma cooperativa de agricultura
familiar. A proposta partiu de um diagnóstico concreto -- informação dispersa,
rastreabilidade frágil, comprovação de entrega precária e baixa visibilidade
gerencial -- e resultou em uma plataforma que conecta produtores, compradores,
motoristas e gestores em um único fluxo digital.

O objetivo geral foi alcançado. A plataforma AgroConecta reúne aplicação Web
administrativa, API REST versionada, aplicativo mobile com operação offline e
banco de dados relacional com procedimentos armazenados e gatilhos. O fluxo
completo foi validado ponta a ponta, em ambiente real: cadastro de produto,
registro de lote com código de rastreio, criação de pedido com itens vinculados
ao lote, programação de rota com geração automática de entrega, atualização de
status pelo motorista, envio de evidência com fotografia e coordenadas, e
emissão de fatura com numeração sequencial e comportamento idempotente.

Os objetivos específicos também foram atendidos. A organização foi caracterizada
com processos, problemas e análise SWOT; a solução foi planejada com personas,
proposta de valor, Lean Canvas e estimativa de viabilidade; as dimensões de
responsabilidade social, diversidade, acessibilidade e proteção de dados foram
incorporadas ao produto, e não apenas ao discurso; a aplicação Web cobre treze
telas com controle de acesso por papel; o aplicativo mobile implementa dez telas
com fila de sincronização persistente e tratamento explícito de conflitos; a
arquitetura em camadas mantém as regras de negócio isoladas de detalhes
tecnológicos; o banco de dados foi modelado em terceira forma normal, com
restrições de integridade, três procedimentos e três gatilhos; a infraestrutura
foi planejada para nuvem, com contêineres, pipeline de integração e entrega
contínuas; e o trabalho foi conduzido com Product Backlog, quatro sprints, quadro
Kanban e critérios de pronto e concluído.

Os resultados quantitativos reforçam a consistência da entrega: 159 testes
automatizados aprovados no backend, com cobertura de 92,5\% no domínio e 93,8\%
na aplicação; 17 suítes e 79 testes no aplicativo mobile; contrato OpenAPI com 22
caminhos; e banco validado em contêiner SQL Server 2022, tanto pelo script
completo quanto pelas migrations. A execução do fluxo real produziu uma fatura
de demonstração e confirmou a persistência de todos os registros.

Do ponto de vista social, a contribuição mais relevante não está na tecnologia em
si, mas na possibilidade de tornar visível a produção de comunidades
tradicionais. Ao registrar a origem do lote, a comunidade do produtor e a
evidência de entrega, a plataforma cria informação que pode ser usada para
acessar políticas públicas, comprovar contrapartidas e reduzir a assimetria de
poder na negociação. A operação offline é parte dessa escolha: sem ela, a
ferramenta excluiria justamente quem mais precisa de apoio.

O trabalho apresenta limitações que devem ser reconhecidas. O APK validado
contempla apenas a arquitetura \texttt{arm64-v8a}; a verificação com leitor de
tela foi parcial; a distribuição em nuvem pelo EAS não foi executada; e o build
da imagem Docker da aplicação não pôde ser concluído no ambiente de
desenvolvimento por restrição de espaço em disco, embora o arquivo esteja
versionado e o pipeline o execute. Tais limitações não invalidam os resultados
funcionais, mas delimitam o alcance da validação.

Como trabalhos futuros, sugere-se: integrar a plataforma a serviços de emissão
de nota fiscal eletrônica; implementar assinatura digital da evidência de
entrega; ampliar o aplicativo para outras arquiteturas e para iOS; incluir
módulo de previsão de demanda a partir do histórico de pedidos; e conduzir um
piloto com produtores reais, medindo tempo de registro, redução de perdas e
satisfação de uso. Essas evoluções preservam a arquitetura atual e aprofundam o
impacto social pretendido.

Conclui-se que o AgroConecta demonstra a aplicação prática e integrada dos
conteúdos do quarto semestre: empreendedorismo, desenvolvimento Web e mobile,
arquitetura de software, programação orientada a objetos, banco de dados,
computação em nuvem, DevOps, gestão ágil e responsabilidade social. Mais do que
um conjunto de artefatos técnicos, o projeto evidencia que a transformação
digital de uma organização depende da articulação entre tecnologia, processo e
pessoas -- e que essa articulação é possível mesmo em contextos de recursos
limitados.
